\documentclass[11pt]{article}
\usepackage[a4paper,textwidth=12.2cm,textheight=21.5cm]{geometry}
\usepackage{fontspec}
\usepackage{setspace}
\setstretch{1.28}
\setmainfont{Times New Roman}
\usepackage{microtype}
\setlength{\emergencystretch}{2em}
\setlength{\parindent}{1.4em}
\setlength{\parskip}{0pt}
\pagestyle{plain}
\begin{document}
\begin{center}
{\LARGE\bfseries Thursdays}
\end{center}
\bigskip
\noindent Kamen Brod has had, in living memory, two barbers. They were brothers. Their shops faced each other across the market street, close enough that each could hear the other's bell, and for fifty years they did not speak.

They had trained at their father's chair, one soaping the faces the other would learn to shave. The father died in the spring of a good year. During the forty days of mourning, the brothers quarreled. What the quarrel was about, the town knew for a while and then lost. The brothers may have kept it. Nobody asked them. Before the forty days were out, Slavcho had taken the corner premises opposite and painted his sign.

The shops were the same shop. Two chairs, a shelf, a list of prices. When Kosta raised the price of a shave, Slavcho's price rose the same week, by no channel anyone could name. They pulled teeth, as barbers did then, and it was said you could tell, in a laugh at the tavern, who went to which.

The town divided as towns do: a man went to the barber his father had gone to, and sent his sons there. Nobody chose. It had been settled for everyone, once, in the forty days. Thursdays, the market day, had always been Slavcho's: the villages came to him, his corner being the first they reached from the bridge.

They opened at the same minute, each blind going up as the other went up, and closed at the same minute, and it could not be said who was following whom. On Sundays they walked to the same church on opposite sides of the street and sat level with each other across the aisle. At funerals — barbers attend funerals; it is trade — they stood at opposite corners of the grave, where the handles would be.

One October they appeared in identical hats, bought the same market-week from different stalls. Neither ever wore his again.

In Kosta's shop, on the nail by the mirror, hung a strop that was not Kosta's. It had been Slavcho's, in their father's time. It was kept oiled. Customers who asked about it once did not ask twice; there was a way he had of stropping his razor, just then, that closed the subject.

Each shaved half the town, and heard, from under his own hands, the other half's news: what Slavcho had said about the new priest, what Kosta had paid for his roof. Each knew the other's life a week late and wrong in the usual places, which is the way the town knew anything.

The new priest, early in his appointment, undertook to reconcile them. He went to each separately and found both men agreeable, and could not afterward say what had gone wrong: each had agreed to everything, contingent on the other's beginning. The priest was new. In time he was not.

When the mill caught fire, the bucket line ran from the river up Well Lane, and the brothers stood in it side by side from midnight until it was out, passing water through four hands as through two. In the morning they were seen to nod. The town lived on it for a year.

Both danced at their niece's wedding, in turn, correctly.

The winter Kosta broke his wrist on the ice by his own step, his shop closed, and his half of the town declined, to a man, to cross the street. For six weeks Kamen Brod went about one-half shaggy, like a hedge cut by a lazy man, and was proud of itself. Then it was found that Slavcho's shop had lately begun opening at dawn, and that men of Kosta's half were coming out of it shaved before the market woke. Nothing was posted. Nothing was charged, or nothing was said to be. When the wrist healed, the dawn stopped.

They grew old at the same rate.

Slavcho died on a Monday morning, in his own chair, between customers. The laying-out was at the house. Kosta came with his case. He shaved his brother with his own hands, which had not touched him since the forty days, and he used the strop from the nail by his mirror. It stropped as if used daily. It had never been used.

Those who closed the shop across the street found, on the nail by Slavcho's mirror, a strop that was not Slavcho's. It was oiled.

Both signs are up. Kosta's shop keeps its days. On Thursdays, he crosses the street.
\end{document}
